\documentclass[aspectratio=169]{beamer}

\usetheme{Madrid}
\usecolortheme{seahorse}

\usepackage[T1]{fontenc}
\usepackage[utf8]{inputenc}
\usepackage[french]{babel}
\usepackage{lmodern}
\usepackage{booktabs}
\usepackage{amsmath}
\usepackage{siunitx}

\title{Synthese Statistique\\Reseaux d'Acteurs IA}
\author{Projet Reseaux d'acteurs}
\date{\today}

\begin{document}

\begin{frame}
  \titlepage
\end{frame}

\begin{frame}{Plan}
\tableofcontents
\end{frame}

\section{Contexte}

\begin{frame}{Perimetre de l'analyse}
\begin{itemize}
  \item Source: API locale et notebook d'analyse.
  \item Population: 1553 entreprises IA.
  \item Couverture geographique: 26 pays.
  \item Granularite sectorielle: 238 secteurs.
  \item Objectif: etablir une vue macro fiable pour prioriser l'enrichissement.
\end{itemize}
\end{frame}

\begin{frame}{Cadre methodologique}
\begin{itemize}
  \item Recuperation de donnees depuis l'API avec pagination et segments (pending, validated).
  \item Uniformisation de la capitalisation en millions (format numerique exploitable).
  \item Production d'indicateurs de completude, validation, geographie et temporalite.
  \item Exports CSV pour reutilisation dans les workflows data et reporting.
\end{itemize}
\end{frame}

\section{Resultats globaux}

\begin{frame}{Indicateurs globaux}
\begin{table}
\centering
\begin{tabular}{lr}
\toprule
Indicateur & Valeur \\
\midrule
Nombre d'entreprises & 1553 \\
Nombre de pays & 26 \\
Nombre de secteurs & 238 \\
Periode couverte & 1975 -- 2024 \\
\bottomrule
\end{tabular}
\end{table}
\end{frame}

\begin{frame}{Lecture macro}
\begin{itemize}
  \item Le corpus est volumineux, mais tres concentre geographiquement.
  \item La couverture sectorielle est riche (238 secteurs), avec un noyau de categories dominantes.
  \item Le signal temporel est recent (mediane 2020) sur les fiches avec annee renseignee.
  \item Le principal goulot de qualite se situe sur les donnees financieres structurees.
\end{itemize}
\end{frame}

\section{Geographie et secteurs}

\begin{frame}{Top 5 des pays}
\begin{table}
\centering
\begin{tabular}{lrr}
\toprule
Pays & Entreprises & Part \\
\midrule
Allemagne & 925 & 59.6\% \\
Suede & 212 & 13.7\% \\
Pays-Bas & 169 & 10.9\% \\
Etats-Unis & 135 & 8.7\% \\
Chine & 16 & 1.0\% \\
\bottomrule
\end{tabular}
\end{table}
\end{frame}

\begin{frame}{Concentration geographique}
\begin{itemize}
  \item L'Allemagne represente la majorite des fiches (59.6\%).
  \item Les trois premiers pays (Allemagne, Suede, Pays-Bas) concentrent 84.2\% du corpus.
  \item Les Etats-Unis et la Chine sont presents mais nettement moins representes.
  \item Implication: les tendances globales refletent fortement la structure du marche allemand.
\end{itemize}
\end{frame}

\begin{frame}{Top 5 des secteurs}
\begin{table}
\centering
\begin{tabular}{lr}
\toprule
Secteur & Entreprises \\
\midrule
AI startup Europe & 193 \\
Human Health \& Social Work Activities & 74 \\
Entreprise IA cotee (market cap) & 61 \\
Cross Industry - Operations & 38 \\
AI startup Europe (Dealroom) & 38 \\
\bottomrule
\end{tabular}
\end{table}
\end{frame}

\begin{frame}{Lecture sectorielle}
\begin{itemize}
  \item Le champ \textit{sector} combine des etiquettes metier et des etiquettes de source.
  \item Certaines categories sont transverses (Cross Industry), d'autres plus verticales (sante).
  \item Une harmonisation semantique des secteurs augmenterait la precision analytique.
  \item Priorite: normaliser les libelles avant comparaison fine inter-secteurs.
\end{itemize}
\end{frame}

\section{Temporalite}

\begin{frame}{Dynamique temporelle}
\begin{itemize}
  \item Annee moyenne de creation: 2017.5
  \item Mediane: 2020
  \item Ecart-type: 8.6
  \item Creations depuis 2020: 47
  \item Part des creations depuis 2020: 54.7\% (sur les fiches avec annee renseignee)
\end{itemize}
\end{frame}

\begin{frame}{Interpretation temporelle}
\begin{itemize}
  \item Le corpus exploitable en chronologie est encore limite (faible taux de fiches datees).
  \item Malgre cela, les donnees disponibles indiquent une acceleration recente.
  \item La dispersion (ecart-type 8.6) suggere un melange d'acteurs etablis et emergents.
  \item Axe d'amelioration: augmenter le remplissage de \textit{founded\_year}.
\end{itemize}
\end{frame}

\section{Qualite des donnees}

\begin{frame}{Qualite et validation des fiches}
\begin{columns}[T]
\column{0.48\textwidth}
\textbf{Validation}
\begin{itemize}
  \item Fiches validees: 22 (1.4\%)
  \item Fiches non validees: 1531 (98.6\%)
\end{itemize}

\column{0.48\textwidth}
\textbf{Completude}
\begin{itemize}
  \item Description: 88.6\%
  \item Website: 89.4\%
  \item Capitalisation renseign\'ee (>0): 9.9\%
\end{itemize}
\end{columns}
\end{frame}

\begin{frame}{Capitalisation (en millions)}
\begin{itemize}
  \item Fiches avec capitalisation > 0: 154 (9.9\%)
  \item Somme des capitalisations: 20\,948\,446.0 M
  \item Mediane: 0.0 M
\end{itemize}

\vspace{0.5em}
\textit{Lecture: une mediane a 0 confirme qu'une majorite de fiches n'a pas encore de capitalisation exploitable.}
\end{frame}

\begin{frame}{Completude: points forts et faibles}
\begin{itemize}
  \item Champs bien renseignes: \textit{name}, \textit{country}, \textit{description}, \textit{website}.
  \item Champs fragiles: \textit{founded\_year}, \textit{employees\_count}, \textit{funds\_raised}.
  \item Donnees financieres encore tres lacunaires malgre la normalisation de format.
  \item La qualite technique est en hausse; l'enjeu restant est la profondeur documentaire.
\end{itemize}
\end{frame}

\begin{frame}{Validation: risque et opportunite}
\begin{itemize}
  \item 1.4\% de fiches validees uniquement.
  \item Risque: analyses decisionnelles biais\'ees par des fiches non confirmees.
  \item Opportunite: industrialiser la revue de fiches par lot (pays/secteur/importance).
  \item Recommendation: introduire un pipeline de validation en trois niveaux (auto, revue, approbation).
\end{itemize}
\end{frame}

\section{Plan d'action}

\begin{frame}{Axes de priorisation}
\begin{enumerate}
  \item Augmenter le taux de validation des fiches critiques (top pays / top secteurs).
  \item Completer en priorite les champs financiers (capitalisation, levees, investisseurs).
  \item Standardiser les metadonnees temporelles pour renforcer les analyses de tendance.
  \item Instrumenter un suivi mensuel: croissance du corpus, validation, completude.
\end{enumerate}
\end{frame}

\begin{frame}{Roadmap operationnelle (90 jours)}
\begin{enumerate}
  \item Semaines 1--3: nettoyage des libelles secteur et revue des pays dominants.
  \item Semaines 4--6: campagne de validation ciblee sur les 200 fiches les plus strategiques.
  \item Semaines 7--9: enrichissement financier prioritaire et consolidation des sources.
  \item Semaines 10--12: recalcul KPI, comparaison baseline, publication d'un rapport executif.
\end{enumerate}
\end{frame}

\begin{frame}{KPI de suivi recommandes}
\begin{table}
\centering
\begin{tabular}{lll}
\toprule
KPI & Baseline & Cible 90j \\
\midrule
Taux de fiches validees & 1.4\% & 15\% \\
Fiches avec capitalisation > 0 & 9.9\% & 25\% \\
Fiches avec annee fondee & 5.5\% & 20\% \\
Fiches avec investisseurs & 1.1\% & 12\% \\
\bottomrule
\end{tabular}
\end{table}
\end{frame}

\begin{frame}{Conclusion}
\begin{itemize}
  \item La base est suffisamment large pour produire des signaux strat\'egiques.
  \item La prochaine marche de valeur passe par la validation et l'enrichissement financier.
  \item Le dispositif analytics est en place (notebook + exports) pour piloter la progression.
  \item La priorisation proposee permet un gain de fiabilite mesurable a court terme.
\end{itemize}
\end{frame}

\begin{frame}
\centering
\Large Merci

\vspace{0.5em}
\normalsize
Questions / Prochaines etapes
\end{frame}

\end{document}
